% Inline Autogyro — Full Assembly Cross-Section v2
% Dual-molding sandwich bearing + swashplate + empennage
% Compile: pdflatex assembly_v2.tex
\documentclass[tikz,border=16pt]{standalone}
\usepackage{amsmath,amssymb,tikz}
\usetikzlibrary{arrows.meta,calc,patterns,decorations.pathreplacing,positioning,shapes.geometric}

\begin{document}
\begin{tikzpicture}[
    scale=1.0,
    dyneema/.style={draw=blue!50!black, line width=2.5pt},
    tube/.style={draw=gray!45, fill=gray!10, line width=1.8pt},
    molding/.style={draw=brown!70!black, fill=brown!15, line width=1.5pt},
    bearing/.style={draw=gray!50, fill=gray!18, line width=1.2pt},
    hub/.style={draw=gray!35, fill=gray!20, line width=1.3pt},
    rotor/.style={draw=red!60!black, fill=red!8, line width=1.6pt},
    blade/.style={draw=red!45!black, fill=red!5, line width=1.4pt},
    actuator/.style={draw=gray!50, fill=gray!25, line width=1.2pt},
    swash/.style={draw=teal!60!black, fill=teal!8, line width=1.3pt},
    empennage/.style={draw=gray!40, fill=gray!18, line width=1.2pt},
    force/.style={-{Stealth[scale=2.0]}, line width=2.2pt},
    forcecomp/.style={-{Stealth[scale=1.5]}, line width=1.6pt, densely dashed},
    annot/.style={font=\footnotesize\sffamily, fill=white, fill opacity=0.90,
                  text opacity=1, inner sep=2pt, rounded corners=1.5pt},
    label/.style={font=\normalsize\sffamily\bfseries},
]

\def\tilt{12}
\def\hubtop{8.8}
\def\hubmid{7.8}
\def\rotory{8.1}
\def\swashy{5.5}
\def\acty{4.2}
\def\empy{2.5}
\def\tubebot{1.5}
\def\splicey{1.5}

% =============================================================
% DYNEEMA LINE
% =============================================================
\draw[dyneema] (0, -0.3) -- (0, 10.3);

% =============================================================
% TUBE (full height, Dyneema through center bore)
% =============================================================
\fill[tube] (-0.35, \tubebot) rectangle (0.35, \hubtop);
\fill[white] (-0.1, \tubebot) rectangle (0.1, \hubtop);
\draw[thin] (-0.35, \tubebot) -- (-0.35, \hubtop);
\draw[thin] (0.35, \tubebot) -- (0.35, \hubtop);

\node[annot] at (0.7, 5.5) {Tube\\\scriptsize(compression\\\scriptsize column)};

% =============================================================
% TOP MOLDING (angled underside, clamped to tube)
% =============================================================
\fill[molding] (-0.5, \hubtop-0.1) rectangle (0.5, \hubtop+0.6);
% Angled cut on underside at tilt
\fill[white, rotate around={-\tilt:(0,\hubtop)}] (-0.6, \hubtop-0.05) rectangle (0.6, \hubtop+0.1);
\fill[molding, rotate around={-\tilt:(0,\hubtop)}] (-0.45, \hubtop-0.05) rectangle (0.45, \hubtop+0.05);
% Clamp bolts
\fill[gray!40] (0.4, \hubtop+0.1) rectangle (0.7, \hubtop+0.35);
\fill[gray!40] (-0.4, \hubtop+0.1) rectangle (-0.7, \hubtop+0.35);

\node[annot] at (1.3, \hubtop+0.3) {Top molding\\(fixed to tube)};

% =============================================================
% TOP THRUST BEARING (under top molding, angled)
% =============================================================
\fill[bearing, rotate around={-\tilt:(0,\hubtop-0.1)}]
    (-0.5, \hubtop-0.35) rectangle (0.5, \hubtop-0.1);
\foreach \x in {-0.35,-0.15,0,0.15,0.35} {
    \fill[gray!55, rotate around={-\tilt:(0,\hubtop-0.22)}]
        (\x, \hubtop-0.25) circle (0.045);
}

\node[annot, right=0.1cm] at (0.8, \hubtop-0.4) {Top thrust\\bearing};

% =============================================================
% ROTOR HUB (hollow, between two bearings, autorotates)
% =============================================================
\fill[hub] (-0.45, \hubmid) rectangle (0.45, \hubtop-0.35);
\fill[white] (-0.12, \hubmid) rectangle (0.12, \hubtop-0.3);
% Angled top face (mates with top bearing)
\fill[hub, rotate around={-\tilt:(0,\hubtop-0.35)}]
    (-0.48, \hubtop-0.35) rectangle (0.48, \hubtop-0.3);
% Angled bottom face
\fill[hub, rotate around={-\tilt:(0,\hubmid)}]
    (-0.48, \hubmid-0.05) rectangle (0.48, \hubmid+0.05);

\node[annot, right=0.2cm] at (0.7, \rotory) {Hollow hub\\(autorotates)};

% =============================================================
% ROTOR DISK + BLADES
% =============================================================
\fill[rotor, rotate around={-\tilt:(0,\rotory)}]
    (-2.5, \rotory-0.12) rectangle (2.5, \rotory+0.12);
\fill[white, rotate around={-\tilt:(0,\rotory)}]
    (-0.5, \rotory-0.15) rectangle (0.5, \rotory+0.15);

% Blades
\fill[blade, rotate around={-\tilt:(0,\rotory)}]
    (-3.5, \rotory-0.12) rectangle (-2.5, \rotory+0.12);
\fill[blade, rotate around={-\tilt:(0,\rotory)}]
    (2.5, \rotory-0.12) rectangle (3.5, \rotory+0.12);

\node[annot] at (-4.2, \rotory) {Blade};

% =============================================================
% BOTTOM THRUST BEARING (under rotor hub)
% =============================================================
\fill[bearing, rotate around={-\tilt:(0,\hubmid)}]
    (-0.5, \hubmid-0.25) rectangle (0.5, \hubmid);
\foreach \x in {-0.35,-0.15,0,0.15,0.35} {
    \fill[gray!55, rotate around={-\tilt:(0,\hubmid-0.12)}]
        (\x, \hubmid-0.15) circle (0.045);
}

% =============================================================
% BOTTOM MOLDING (angled topside, clamped to tube)
% =============================================================
\fill[molding] (-0.5, \hubmid-0.55) rectangle (0.5, \hubmid-0.25);
\fill[white, rotate around={-\tilt:(0,\hubmid)}] (-0.6, \hubmid-0.35) rectangle (0.6, \hubmid);
\fill[molding, rotate around={-\tilt:(0,\hubmid-0.25)}]
    (-0.48, \hubmid-0.3) rectangle (0.48, \hubmid-0.2);

\node[annot, right=0.1cm] at (0.8, \hubmid-0.7) {Bottom molding\\(clamped to tube)};

% =============================================================
% SANDWICH CLAMP ANNOTATION
% =============================================================
\draw[forcecomp, brown!60!black] (0, \hubtop+0.5) -- (0, \hubtop);
\draw[forcecomp, brown!60!black] (0, \hubmid-0.6) -- (0, \hubmid);
\node[annot, brown!60!black] at (-1.4, \rotory)
    {Moldings clamp\\rotor hub via\\bearings above\\and below};

% =============================================================
% SWASHPLATE (below rotor, around tube)
% =============================================================
% Rotating ring (upper)
\fill[swash] (-0.5, \swashy+0.3) rectangle (0.5, \swashy+0.55);
% Stationary ring (lower)
\fill[swash!70!gray] (-0.55, \swashy) rectangle (0.55, \swashy+0.25);

\node[annot, right=0.2cm] at (0.8, \swashy+0.4) {Swashplate\\(rotating ring)};
\node[annot, right=0.2cm] at (0.8, \swashy+0.1) {Stationary ring};

% Scissor link (hub to rotating ring)
\draw[forcecomp, gray] (0.45, \hubmid) -- (0.5, \swashy+0.55);
\draw[forcecomp, gray] (-0.45, \hubmid) -- (-0.5, \swashy+0.55);

% =============================================================
% ACTUATORS (3× around tube, pushrods to stationary ring)
% =============================================================
\fill[actuator] (-0.25, \acty) rectangle (0.05, \acty+1.0);
\fill[actuator] (0.2, \acty) rectangle (0.5, \acty+1.0);
\draw[forcecomp, gray] (0.35, \acty+1.0) -- (0.45, \swashy);
\draw[forcecomp, gray] (-0.1, \acty+1.0) -- (-0.35, \swashy);

\node[annot, right=0.2cm] at (0.7, \acty+0.6) {3× actuators\\(linear servos)};

% =============================================================
% EMPENNAGE (bottom of tube, extends downwind)
% =============================================================
% Tail boom bracket
\fill[empennage] (0.3, \empy-0.15) rectangle (0.5, \empy+0.15);
% Tail boom
\fill[empennage] (0.4, \empy-0.1) rectangle (2.5, \empy+0.1);
% H-stab
\fill[empennage] (2.0, \empy-0.1) rectangle (2.5, 3.6);
% V-fin (hangs below boom, downwind side)
\fill[empennage] (2.0, 2.5) rectangle (2.3, 1.5);
\node[annot] at (2.8, 2.0) {V-fin};
\node[annot] at (3.3, \empy) {H-stab};
\node[annot] at (1.8, \empy+0.9) {Tail boom};

% =============================================================
% WEBBING CAPTURE (bottom of tube → spliced eye)
% =============================================================
% Splice in Dyneema
\fill[brown!30] (-0.15, \tubebot-0.5) rectangle (0.15, \tubebot+0.2);
% Webbing loops
\draw[orange!70!black, line width=2.0pt]
    (0, \tubebot) -- (1.5, \tubebot) -- (1.5, \tubebot-0.8) -- (0.35, \tubebot-0.8);
\draw[orange!70!black, line width=2.0pt]
    (0, \tubebot) -- (-1.5, \tubebot) -- (-1.5, \tubebot-0.8) -- (-0.35, \tubebot-0.8);

\node[annot, orange!70!black] at (2.5, \tubebot-0.3) {Spectra webbing\\(tube capture)};
\node[annot, brown!70] at (-1.6, \tubebot-0.2) {Spliced eye\\(in Dyneema)};

% =============================================================
% FORCE PATHS
% =============================================================
% Lift from rotor
\draw[force, green!50!black, very thick]
    (0, \rotory+0.5) -- ++({1.8*sin(\tilt)}, {1.8*cos(\tilt)})
    node[right, font=\footnotesize, green!50!black] {$F_\text{lift}$};

% Compression through tube
\draw[force, purple!50!black, line width=3.5pt]
    (0, \hubtop-0.3) -- (0, \tubebot);
\draw[force, purple!50!black, line width=3.5pt]
    (0, \hubmid-0.5) -- (0, \tubebot);

% Tension above/below
\node[annot, purple!50!black, left=0.1cm] at (-0.5, 9.8) {$T_\text{above}$};
\node[annot, purple!50!black, left=0.1cm] at (-0.5, 0.2) {$T_\text{below}$};
\draw[force, purple!50!black] (-0.4, 9.5) -- (-0.4, \hubtop+0.5);
\draw[force, purple!50!black] (0.4, \tubebot-0.8) -- (0.4, 0.3);

% =============================================================
% LABELS & NOTES
% =============================================================
\node[draw=black!25, fill=white, rounded corners=3pt, inner sep=7pt,
      anchor=north west, font=\scriptsize\sffamily, align=left, text width=5cm]
    at (5.0, 9.8) {
    \textbf{Assembly order (top → bottom)}\\[3pt]
    1. Top molding (angled underside)\\
    2. Top thrust bearing\\
    3. Rotor hub + disk + blades\\
    \quad (autorotates between bearings)\\
    4. Bottom thrust bearing\\
    5. Bottom molding\\
    6. Swashplate (below rotor)\\
    7. 3× linear actuators\\
    8. Empennage boom + H-stab + V-fin\\
    9. Webbing capture → spliced eye\\[4pt]
    \textbf{Load path:}\\
    $F_\text{lift}$ → hub → top bearing →\\
    top molding → tube compression →\\
    bottom molding → webbing → Dyneema\\
    \smallskip
    $T_\text{below} = T_\text{above} + F_\text{lift} - W$
};

% Title
\node[label, font=\Large] at (5.5, 10.8)
    {Inline Autogyro Rotor Assembly — Full Cross-Section};

\end{tikzpicture}
\end{document}
